\subsection*{CU-08: Configurar la cuenta la primera vez}
\addcontentsline{toc}{subsection}{CU-08: Configurar la cuenta la primera vez}
\label{cu-08}

\begin{fichacu}{CU-08}{Configurar la cuenta la primera vez}
  \crudFila{Responsable}{José Eduardo Prior Hernández}
  \crudFila{Actor principal}{Jugador o Invitado}
  \crudFila{Actores de sistema}{Servidor de partidas, Base de datos}
  \crudFila{Prototipos}{6d, 4a}
  \crudFila{Objetivo}{Dar al jugador, en su primer acceso, la ocasión de elegir el peón con el que va a jugar y el icono con el que lo verán los demás, en un solo paso y sin obligarlo a recorrer la pantalla de personalización completa.}
  \crudFila{Disparador}{El SJM recibe el perfil de una cuenta cuya \texttt{fecha\_\allowbreak{}configuracion\_\allowbreak{}inicial} es nula, al terminar el \mbox{CU-01} por su \mbox{FA-18} o al terminar el \mbox{CU-06}.}
  \textbf{Precondiciones} & \cuCelda{\begin{cureglas}
      \item \textbf{\mbox{PRE-01}} El jugador tiene una \textit{Sesion} abierta y su token es válido.
      \item \textbf{\mbox{PRE-02}} El \textit{Usuario} tiene \texttt{fecha\_\allowbreak{}configuracion\_\allowbreak{}inicial} nula.
      \item \textbf{\mbox{PRE-03}} El \textit{Usuario} posee, mediante \textit{UsuarioObjeto}, al menos un \textit{ObjetoCosmetico} inicial de la \textit{Ranura} de peón.
      \item \textbf{\mbox{PRE-04}} El catálogo de iconos predefinidos contiene los treinta y dos iconos.
      \item \textbf{\mbox{PRE-05}} El Servidor de partidas está en ejecución y tiene conexión disponible con la Base de datos.
    \end{cureglas}} \\ \hline
  \textbf{Flujo normal} & \cuCelda{\begin{cuenum}[start=1]
      \item El SJM despliega la \texttt{GUI\_\allowbreak{}FirstTime}, que muestra en un solo paso los peones iniciales que posee el \textit{Usuario} sobre un tablero de vista previa, los treinta y dos iconos predefinidos, el \texttt{nickname} de la cuenta y las opciones ``Empezar'' y ``Omitir''. El peón y el icono predeterminados aparecen ya seleccionados. \textit{(Ver \mbox{FA-01})}
      \item El jugador selecciona un peón. \textit{(Ver \mbox{FA-02})}
      \item El SJM coloca el peón elegido en la vista previa, sin haberlo equipado todavía.
      \item El jugador selecciona un icono.
      \item El jugador da click en ``Empezar''.
      \item El SJM envía el mensaje \texttt{FIRST\_\allowbreak{}TIME\_\allowbreak{}SETUP\_\allowbreak{}REQUEST} con el \texttt{id\_\allowbreak{}objeto} del peón, el \texttt{id\_\allowbreak{}icono} y el token de sesión. \textit{(Ver \mbox{EX-03}, \mbox{EX-04}, \mbox{EX-05})}
      \item El Servidor de partidas verifica que el token corresponda a una \textit{Sesion} abierta. \textit{(Ver \mbox{FA-03})}
      \item El Servidor de partidas verifica que \texttt{fecha\_\allowbreak{}configuracion\_\allowbreak{}inicial} del \textit{Usuario} siga nula. \textit{(Ver \mbox{FA-04})}
    \end{cuenum}} \\
   & \cuCelda{\begin{cuenum}[start=9]
      \item El Servidor de partidas verifica que el peón recibido sea un \textit{ObjetoCosmetico} de la \textit{Ranura} de peón que el \textit{Usuario} posee, es decir, que exista su \textit{UsuarioObjeto} (\mbox{RN-02}). \textit{(Ver \mbox{FA-05})}
      \item El Servidor de partidas verifica que el \texttt{id\_\allowbreak{}icono} esté entre los treinta y dos iconos predefinidos (\mbox{RN-03}). \textit{(Ver \mbox{FA-05})}
      \item El Servidor de partidas inicia una transacción sobre la Base de datos.
      \item El Servidor de partidas actualiza la fila de \textit{Equipamiento} de la \textit{Ranura} de peón para que apunte al peón elegido. \textit{(Ver \mbox{EX-01})}
      \item El Servidor de partidas actualiza en el \textit{Usuario} el \texttt{id\_\allowbreak{}icono} y escribe \texttt{fecha\_\allowbreak{}configuracion\_\allowbreak{}inicial} con la fecha y hora actuales (\mbox{RN-01}).
      \item El Servidor de partidas confirma la transacción. \textit{(Ver \mbox{EX-02})}
      \item El Servidor de partidas responde con \texttt{FIRST\_\allowbreak{}TIME\_\allowbreak{}SETUP\_\allowbreak{}OK} y el equipamiento actualizado.
      \item El SJM despliega la \texttt{GUI\_\allowbreak{}MainMenu} con el peón y el icono elegidos. \textit{(Ver \mbox{FA-06})}
    \end{cuenum}} \\
   & \cuCelda{\begin{cuenum}[start=17]
      \item Fin del caso de uso.
    \end{cuenum}} \\ \hline
  \textbf{Flujos alternos} & \cuCelda{\textbf{\mbox{FA-01}: El jugador omite la configuración}\par
    \begin{cuenum}[start=1]
      \item El jugador selecciona ``Omitir''.
      \item El SJM envía el mensaje \texttt{FIRST\_\allowbreak{}TIME\_\allowbreak{}SETUP\_\allowbreak{}REQUEST} sin peón ni icono.
      \item El Servidor de partidas conserva el peón y el icono predeterminados, y escribe únicamente \texttt{fecha\_\allowbreak{}configuracion\_\allowbreak{}inicial} (\mbox{RN-04}).
      \item El flujo continúa en el paso 15 del FN. La pantalla no volverá a aparecer: el peón y el icono pueden cambiarse después desde el \mbox{CU-14} y el \mbox{CU-12}.
    \end{cuenum}} \\
   & \cuCelda{\textbf{\mbox{FA-02}: El jugador quiere ver el peón sobre el tablero completo}\par
    \begin{cuenum}[start=1]
      \item El jugador arrastra la vista previa para girarla.
      \item El SJM rota el tablero de vista previa con el peón seleccionado, con la misma cámara libre de la partida.
      \item El flujo continúa en el paso 2 del FN.
    \end{cuenum}} \\
   & \cuCelda{\textbf{\mbox{FA-03}: La sesión ya no es válida}\par
    \begin{cuenum}[start=1]
      \item El Servidor de partidas responde con \texttt{SESSION\_\allowbreak{}INVALID}.
      \item El SJM muestra la \texttt{GUI\_\allowbreak{}MessageWarning} indicando que la sesión terminó y despliega la \texttt{GUI\_\allowbreak{}Login}. La configuración inicial queda pendiente para el siguiente acceso.
      \item Fin del caso de uso.
    \end{cuenum}} \\
   & \cuCelda{\textbf{\mbox{FA-04}: La configuración inicial ya se había hecho}\par
    \begin{cuenum}[start=1]
      \item El Servidor de partidas encuentra \texttt{fecha\_\allowbreak{}configuracion\_\allowbreak{}inicial} con valor: la cuenta ya la completó, por ejemplo desde otro dispositivo abierto al mismo tiempo (\mbox{D-01}).
      \item El Servidor de partidas no modifica nada y responde con \texttt{FIRST\_\allowbreak{}TIME\_\allowbreak{}SETUP\_\allowbreak{}ALREADY\_\allowbreak{}DONE} y el equipamiento vigente.
      \item El SJM despliega la \texttt{GUI\_\allowbreak{}MainMenu} con el equipamiento recibido, sin mostrar ningún error.
      \item Fin del caso de uso.
    \end{cuenum}} \\
   & \cuCelda{\textbf{\mbox{FA-05}: El peón o el icono recibidos no son válidos}\par
    \begin{cuenum}[start=1]
      \item El Servidor de partidas registra la discrepancia en la bitácora del servidor con la \texttt{version\_\allowbreak{}cliente} y la dirección de red de origen: la interfaz solo ofrece peones poseídos e iconos del catálogo, así que es indicio de cliente modificado.
      \item El Servidor de partidas responde con \texttt{FIRST\_\allowbreak{}TIME\_\allowbreak{}SETUP\_\allowbreak{}FAILED} y el motivo \texttt{DATOS\_\allowbreak{}INVALIDOS}.
      \item El SJM recarga los peones y los iconos disponibles y muestra la \texttt{GUI\_\allowbreak{}MessageWarning}.
      \item El flujo continúa en el paso 2 del FN.
    \end{cuenum}} \\
   & \cuCelda{\textbf{\mbox{FA-06}: El jugador es un invitado}\par
    \begin{cuenum}[start=1]
      \item El SJM añade a la \texttt{GUI\_\allowbreak{}MainMenu} el aviso permanente de cuenta no vinculada del \mbox{CU-06}.
      \item Fin del caso de uso. La configuración elegida se conserva si más adelante vincula la cuenta con el \mbox{CU-07}.
    \end{cuenum}} \\ \hline
  \textbf{Excepciones} & \cuCelda{\textbf{\mbox{EX-01}: Fallo al escribir en la base de datos}\par
    \begin{cuenum}[start=1]
      \item El Servidor de partidas revierte la transacción, de modo que no quede el peón equipado sin \texttt{fecha\_\allowbreak{}configuracion\_\allowbreak{}inicial}, ni al revés.
      \item El Servidor de partidas registra la excepción en la bitácora del servidor con un identificador de incidencia y responde con \texttt{SERVER\_\allowbreak{}ERROR}.
      \item El SJM muestra la \texttt{GUI\_\allowbreak{}MessageError} con el identificador y conserva la selección.
      \item El flujo continúa en el paso 5 del FN.
    \end{cuenum}} \\
   & \cuCelda{\textbf{\mbox{EX-02}: Fallo al confirmar la transacción}\par
    \begin{cuenum}[start=1]
      \item El Servidor de partidas trata la transacción como fallida y registra la excepción señalando que el estado es indeterminado.
      \item El Servidor de partidas responde con \texttt{SERVER\_\allowbreak{}ERROR}.
      \item El SJM despliega la \texttt{GUI\_\allowbreak{}MainMenu} con el equipamiento predeterminado. Si la configuración no se escribió, la pantalla volverá a aparecer en el siguiente acceso; si sí se escribió, el \mbox{FA-04} lo resolverá sin repetirla.
      \item Fin del caso de uso.
    \end{cuenum}} \\
   & \cuCelda{\textbf{\mbox{EX-03}: Se agota el tiempo de espera de la respuesta}\par
    \begin{cuenum}[start=1]
      \item El SJM detecta que transcurrió el tiempo límite sin recibir un mensaje completo.
      \item El SJM no reenvía la configuración por su cuenta y muestra la \texttt{GUI\_\allowbreak{}MessageError} con las opciones ``Reintentar'' y ``Continuar''.
      \item Si el jugador selecciona ``Reintentar'', el flujo continúa en el paso 6 del FN; un reintento sobre una configuración ya escrita termina en el \mbox{FA-04}. Si selecciona ``Continuar'', el SJM despliega la \texttt{GUI\_\allowbreak{}MainMenu}.
    \end{cuenum}} \\
   & \cuCelda{\textbf{\mbox{EX-04}: La conexión se interrumpe durante el intercambio}\par
    \begin{cuenum}[start=1]
      \item El SJM detecta el cierre inesperado del socket y muestra la barra de conexión perdida.
      \item Al recuperar la conexión, el SJM vuelve a solicitar el perfil y, si \texttt{fecha\_\allowbreak{}configuracion\_\allowbreak{}inicial} sigue nula, el flujo continúa en el paso 1 del FN.
    \end{cuenum}} \\
   & \cuCelda{\textbf{\mbox{EX-05}: El mensaje recibido está malformado}\par
    \begin{cuenum}[start=1]
      \item El SJM descarta el mensaje, cierra la conexión y registra en la bitácora local el contenido truncado.
      \item El SJM muestra la \texttt{GUI\_\allowbreak{}MessageError} indicando que la respuesta no pudo interpretarse.
      \item Fin del caso de uso.
    \end{cuenum}} \\ \hline
  \textbf{Postcondiciones} & \cuCelda{\begin{cureglas}
      \item \textbf{\mbox{POST-01}} Si el caso termina por el FN, la fila de \textit{Equipamiento} de la \textit{Ranura} de peón apunta al peón elegido y las otras cinco no cambiaron.
      \item \textbf{\mbox{POST-02}} Si el caso termina por el FN, el \textit{Usuario} tiene el \texttt{id\_\allowbreak{}icono} elegido.
      \item \textbf{\mbox{POST-03}} Si el caso termina por el FN o por el \mbox{FA-01}, \texttt{fecha\_\allowbreak{}configuracion\_\allowbreak{}inicial} tiene valor y la pantalla no vuelve a aparecer en ningún dispositivo.
      \item \textbf{\mbox{POST-04}} Ningún \textit{UsuarioObjeto} se creó ni se eliminó: elegir el peón inicial no otorga nada que la cuenta no tuviera.
      \item \textbf{\mbox{POST-05}} Si el caso termina por una excepción antes de escribir, \texttt{fecha\_\allowbreak{}configuracion\_\allowbreak{}inicial} sigue nula y la configuración se ofrecerá en el siguiente acceso.
    \end{cureglas}} \\ \hline
  \textbf{Reglas de negocio} & \cuCelda{\begin{cureglas}
      \item \textbf{\mbox{RN-01}} La configuración inicial se ofrece una sola vez por cuenta. La marca es \texttt{fecha\_\allowbreak{}configuracion\_\allowbreak{}inicial}, y no un indicador booleano, porque la fecha sirve además para saber cuánto tardó el jugador en empezar.
      \item \textbf{\mbox{RN-02}} El peón elegido debe ser un objeto inicial de la \textit{Ranura} de peón que el \textit{Usuario} ya posee. Esta pantalla no otorga objetos: el alta ya otorgó todos los iniciales (\mbox{CU-02} paso 18).
      \item \textbf{\mbox{RN-03}} El icono se elige entre los treinta y dos iconos predefinidos de animales y objetos. Subir una imagen propia no está disponible aquí; es el \mbox{FA-01} del \mbox{CU-12}.
      \item \textbf{\mbox{RN-04}} Omitir la configuración es válido y equivale a aceptar los predeterminados. La pantalla no insiste.
      \item \textbf{\mbox{RN-05}} El \texttt{nickname} no se elige en este caso: se eligió en el registro (\mbox{CU-02}) o se elegirá al vincular la cuenta de invitado (\mbox{CU-07}).
      \item \textbf{\mbox{RN-06}} Todas las validaciones del cliente se repiten en el servidor.
    \end{cureglas}} \\ \hline
  \textbf{Observación} & \cuCelda{\textit{Sobre por qué esta pantalla no pide el nombre.}\par
    \texttt{descripcion.md} dice que la primera vez se eligen nombre, peón e icono. El nombre, sin embargo, ya se pide en el registro, porque sirve como identificador de acceso y debe ser único antes de que exista la cuenta. Pedirlo aquí obligaría a crear cuentas sin \texttt{nickname}, lo que rompería la unicidad y el inicio de sesión por nombre. Se conservó la intención del diseño —que el primer contacto sea elegir cómo te ven— reduciendo la pantalla a lo que todavía no se ha decidido: el peón y el icono.} \\ \hline
  \textbf{Datos que toca} & \cuCelda{\begin{cudatos}
      \item \textbf{\textit{Sesion}:} lee por token \textit{(paso 7)}
      \item \textbf{\textit{Usuario}:} lee \texttt{fecha\_\allowbreak{}configuracion\_\allowbreak{}inicial}, \texttt{nickname} \textit{(pasos 1, 8)}
      \item \textbf{\textit{Usuario}:} escribe \texttt{id\_\allowbreak{}icono}, \texttt{fecha\_\allowbreak{}configuracion\_\allowbreak{}inicial} \textit{(pasos 13, \mbox{FA-01})}
      \item \textbf{\textit{UsuarioObjeto}, \textit{ObjetoCosmetico}:} lee los peones iniciales poseídos \textit{(pasos 1, 9)}
      \item \textbf{\textit{Equipamiento}:} actualiza la fila de la ranura de peón \textit{(paso 12)}
    \end{cudatos}} \\ \hline
\end{fichacu}
\clearpage
